\section{Condiciones de frontera y coeficientes de Fresnel en dieléctricos}
\label{sec:appendix_maxwell_frontera_fresnel}

En este apéndice se deducen las condiciones de frontera de Maxwell en una interfaz dieléctrica y las expresiones de reflexión/transmisión para ondas planas en polarizaciones $s$ y $p$.

\subsection{Condiciones de frontera en una interfaz dieléctrica}
\label{subsec:apx_condiciones_frontera}

Partimos de las ecuaciones de Maxwell en medios materiales,
\begin{equation}
\begin{aligned}
\nabla \cdot \mathbf{D} &= \rho_f, & \nabla \cdot \mathbf{B} &= 0, \\
\nabla \times \mathbf{E} &= -\frac{\partial \mathbf{B}}{\partial t},
& \nabla \times \mathbf{H} &= \mathbf{J}_f + \frac{\partial \mathbf{D}}{\partial t}.
\end{aligned}
\label{eq:apx_maxwell_materiales}
\end{equation}

Sea la interfaz
\begin{equation}
\Sigma = \{(x,y,z)\mid z=0\},
\end{equation}
que separa dos medios ($z<0$ y $z>0$). Para un campo vectorial genérico,
\begin{equation}
\tilde{\mathbf F}(\mathbf r,t)
=
\mathbf F(\mathbf r,t)\,\Theta(-z)
+
\mathbf F'(\mathbf r,t)\,\Theta(z),
\label{eq:apx_campo_partido}
\end{equation}
con
\begin{equation}
\Theta(z)=
\begin{cases}
1, & z>0,\\
0, & z\le 0.
\end{cases}
\end{equation}

La divergencia y el rotacional toman la forma distribucional:
\begin{align}
\nabla\cdot\tilde{\mathbf F}
&=
\left[\nabla\cdot\mathbf F\right]_H
+
\delta(z)\,\hat{\mathbf n}\cdot\left(\mathbf F'-\mathbf F\right),
\label{eq:apx_div_partida}\\
\nabla\times\tilde{\mathbf F}
&=
\left[\nabla\times\mathbf F\right]_H
+
\delta(z)\,\hat{\mathbf n}\times\left(\mathbf F'-\mathbf F\right).
\label{eq:apx_rot_partida}
\end{align}

Descomponemos fuentes volumétricas y superficiales como
\begin{equation}
\rho_f = \tilde\rho_f + \sigma_f(x,y,t)\,\delta(z),
\qquad
\mathbf J_f = \tilde{\mathbf J}_f + \mathbf K_f(x,y,t)\,\delta(z).
\label{eq:apx_fuentes_partidas}
\end{equation}

Sustituyendo \eqref{eq:apx_div_partida}, \eqref{eq:apx_rot_partida} y \eqref{eq:apx_fuentes_partidas} en \eqref{eq:apx_maxwell_materiales}, y agrupando términos proporcionales a $\delta(z)$, se obtienen las condiciones de frontera:
\begin{equation}
\begin{aligned}
\hat{\mathbf n}\cdot(\mathbf D_2-\mathbf D_1) &= \sigma_f,\\
\hat{\mathbf n}\cdot(\mathbf B_2-\mathbf B_1) &= 0,\\
\hat{\mathbf n}\times(\mathbf E_2-\mathbf E_1) &= 0,\\
\hat{\mathbf n}\times(\mathbf H_2-\mathbf H_1) &= \mathbf K_f.
\end{aligned}
\label{eq:apx_condiciones_frontera_finales}
\end{equation}

\subsection{Ondas planas y relaciones de dispersión}
\label{subsec:apx_ondas_planas}

En ausencia de fuentes libres ($\rho=0$, $\mathbf J=0$):
\begin{equation}
\begin{aligned}
\nabla \cdot \mathbf E &= 0, & \nabla \cdot \mathbf B &= 0,\\
\nabla \times \mathbf E &= -\frac{\partial \mathbf B}{\partial t},
& \nabla \times \mathbf B &= \frac{1}{c^2}\frac{\partial \mathbf E}{\partial t},
\end{aligned}
\label{eq:apx_maxwell_vacio_local}
\end{equation}
con
\begin{equation}
c = \frac{1}{\sqrt{\mu\epsilon}}.
\label{eq:apx_c_medio}
\end{equation}

Para una onda plana armónica,
\begin{equation}
\mathbf E = \mathbf E_0 e^{i(\mathbf k\cdot\mathbf r-\omega t)},
\qquad
\mathbf B = \mathbf B_0 e^{i(\mathbf k\cdot\mathbf r-\omega t)},
\label{eq:apx_ansatz_onda_plana}
\end{equation}
se obtiene
\begin{equation}
\begin{aligned}
i\,\mathbf k\cdot\mathbf E &= 0,
& i\,\mathbf k\cdot\mathbf B &= 0,\\
i\,\mathbf k\times\mathbf E &= i\omega\mathbf B,
& \mathbf k\times\mathbf B &= -\frac{\omega}{c^2}\mathbf E.
\end{aligned}
\label{eq:apx_maxwell_k_omega}
\end{equation}

Aplicando $\mathbf k\times(\mathbf k\times\mathbf E)=\mathbf k(\mathbf k\cdot\mathbf E)-k^2\mathbf E$ y usando \eqref{eq:apx_maxwell_k_omega}, resulta
\begin{equation}
k^2 = \frac{\omega^2}{c^2},
\qquad
\omega = ck,
\qquad
v = \frac{\omega}{k}.
\label{eq:apx_dispersion}
\end{equation}

\subsection{Interfaz plana y coeficientes de Fresnel}
\label{subsec:apx_fresnel}

Consideramos tres ondas planas en la interfaz $\Sigma$: incidente, reflejada y transmitida,
\begin{equation}
\mathbf{E}_j = \mathbf{E}_j^0\, e^{i(\mathbf{k}_j\cdot\mathbf{r} - \omega t)},
\qquad j\in\{i,R,T\}.
\label{eq:apx_campos_iRT}
\end{equation}

Situamos la interfaz en $z=0$ con normal $\hat{\mathbf{n}} = \hat{\mathbf{z}}$ apuntando hacia el medio incidente. Los vectores de onda unitarios son
\begin{equation}
\hat{\mathbf{k}}_i = \sin\theta_i\,\hat{\mathbf{x}} - \cos\theta_i\,\hat{\mathbf{z}},
\quad
\hat{\mathbf{k}}_R = \sin\theta_i\,\hat{\mathbf{x}} + \cos\theta_i\,\hat{\mathbf{z}},
\quad
\hat{\mathbf{k}}_T = \sin\theta_T\,\hat{\mathbf{x}} - \cos\theta_T\,\hat{\mathbf{z}}.
\label{eq:apx_kvectors}
\end{equation}

El \textit{plano de incidencia} es $\Pi = \operatorname{span}\{\hat{\mathbf{k}}_i, \hat{\mathbf{n}}\}$, que en este caso coincide con el plano $xz$. Se definen dos polarizaciones ortogonales:
\begin{itemize}
\item \textit{Polarización $s$} (\textit{senkrecht}, perpendicular a $\Pi$): $\hat{\mathbf{s}} = \hat{\mathbf{y}}$, común a los tres campos.
\item \textit{Polarización $p$} (\textit{parallel}, en $\Pi$ con $\hat{\mathbf{p}}_j \perp \hat{\mathbf{k}}_j$): se obtiene rotando $\hat{\mathbf{k}}_j$ en $\Pi$ un ángulo $\pi/2$ hacia $\hat{\mathbf{n}}$,
\begin{equation}
\hat{\mathbf{p}}_i = \cos\theta_i\,\hat{\mathbf{x}} + \sin\theta_i\,\hat{\mathbf{z}},
\quad
\hat{\mathbf{p}}_R = -\cos\theta_i\,\hat{\mathbf{x}} + \sin\theta_i\,\hat{\mathbf{z}},
\quad
\hat{\mathbf{p}}_T = \cos\theta_T\,\hat{\mathbf{x}} + \sin\theta_T\,\hat{\mathbf{z}}.
\label{eq:apx_pvectors}
\end{equation}
\end{itemize}

Cada amplitud vectorial se descompone como
\begin{equation}
\mathbf{E}_j^0 = E_j^s\,\hat{\mathbf{s}} + E_j^p\,\hat{\mathbf{p}}_j.
\label{eq:apx_descomposicion_sp}
\end{equation}

Como $s$ y $p$ son ortogonales, las condiciones de frontera \eqref{eq:apx_condiciones_frontera_finales} se desacoplan para cada modo. Definimos también
\begin{equation}
Z = \sqrt{\frac{\epsilon}{\mu}},
\qquad
Z' = \sqrt{\frac{\epsilon'}{\mu'}},
\label{eq:apx_def_Z}
\end{equation}
y los coeficientes $R_j = E_R^j/E_i^j$, $T_j = E_T^j/E_i^j$ para $j \in \{s, p\}$.

\paragraph{Modo $s$.}

La componente $s$ es tangencial a $\Sigma$, por lo que la condición iii) da directamente
\begin{equation}
E_i^s + E_R^s = E_T^s.
\label{eq:apx_cond_s_E}
\end{equation}

Para la condición iv) se aplica la identidad $(\mathbf{k}\times\mathbf{E}^s)\times\hat{\mathbf{n}} = (\hat{\mathbf{k}}\cdot\hat{\mathbf{n}})\mathbf{E}^s$, válida porque $\hat{\mathbf{n}}\cdot\hat{\mathbf{s}} = 0$. Con los productos $\hat{\mathbf{k}}_i\cdot\hat{\mathbf{n}} = -\cos\theta_i$, $\hat{\mathbf{k}}_R\cdot\hat{\mathbf{n}} = \cos\theta_i$, $\hat{\mathbf{k}}_T\cdot\hat{\mathbf{n}} = -\cos\theta_T$ y usando $k/\mu = \omega\sqrt{\epsilon/\mu} = \omega Z$, se obtiene
\begin{equation}
Z\cos\theta_i\,(E_i^s - E_R^s) = Z'\cos\theta_T\,E_T^s.
\label{eq:apx_cond_s_H}
\end{equation}

Resolviendo el sistema \eqref{eq:apx_cond_s_E}--\eqref{eq:apx_cond_s_H}:
\begin{align}
R_s &= \frac{Z\cos\theta_i - Z'\cos\theta_T}{Z\cos\theta_i + Z'\cos\theta_T},
\label{eq:apx_fresnel_rs}\\
T_s &= \frac{2Z\cos\theta_i}{Z\cos\theta_i + Z'\cos\theta_T}.
\label{eq:apx_fresnel_ts}
\end{align}

\paragraph{Modo $p$.}

La componente tangencial ($\hat{\mathbf{x}}$) de la condición iii), usando que $\hat{\mathbf{p}}_R\cdot\hat{\mathbf{x}} = -\cos\theta_i$, da
\begin{equation}
(E_i^p - E_R^p)\cos\theta_i = E_T^p\cos\theta_T.
\label{eq:apx_cond_p_E}
\end{equation}

El campo magnético del modo $p$ es $\mathbf{H}_j = \frac{k_j}{\omega\mu_j}(\hat{\mathbf{k}}_j \times \hat{\mathbf{p}}_j)E_j^p$; de \eqref{eq:apx_kvectors} y \eqref{eq:apx_pvectors} se comprueba que $\hat{\mathbf{k}}_j\times\hat{\mathbf{p}}_j = -\hat{\mathbf{s}}$ para los tres campos, de modo que $\mathbf{H}_j^p$ es tangencial a $\Sigma$ con magnitud $Z_j E_j^p$. La condición iv) da
\begin{equation}
Z\,(E_i^p + E_R^p) = Z'\,E_T^p.
\label{eq:apx_cond_p_H}
\end{equation}

Resolviendo el sistema \eqref{eq:apx_cond_p_E}--\eqref{eq:apx_cond_p_H}:
\begin{align}
R_p &= \frac{Z'\cos\theta_i - Z\cos\theta_T}{Z'\cos\theta_i + Z\cos\theta_T},
\label{eq:apx_fresnel_rp}\\
T_p &= \frac{2Z\cos\theta_i}{Z'\cos\theta_i + Z\cos\theta_T}.
\label{eq:apx_fresnel_tp}
\end{align}

\paragraph{Expresión en términos de índices de refracción y campo transmitido.}

Para medios no magnéticos ($\mu = \mu' = \mu_0$) se tiene $Z_j = n_j/Z_0$ con $Z_0 = \sqrt{\mu_0/\epsilon_0}$, y los factores $Z_0$ se cancelan. Los coeficientes de transmisión quedan
\begin{align}
t_s &= \frac{2n_1\cos\theta_i}{n_1\cos\theta_i + n_2\cos\theta_t},
\label{eq:apx_fresnel_ts_n}\\
t_p &= \frac{2n_1\cos\theta_i}{n_2\cos\theta_i + n_1\cos\theta_t},
\label{eq:apx_fresnel_tp_n}
\end{align}
y el campo eléctrico transmitido sobre la interfaz resulta
\begin{equation}
\vb{E}'_\Sigma
= t_p\,(\vb{E}_\Sigma\cdot\uvec{p})\,\uvec{p}'
+ t_s\,(\vb{E}_\Sigma\cdot\uvec{s})\,\uvec{s}'.
\label{eq:apx_E_transmitido_Sigma}
\end{equation}
